% !TEX TS-program = lualatex+se+debug
% !TEX encoding = UTF-8 Unicode

\documentclass[14pt, twoside, a4paper]{extreport}
% -----------------------------------------------------------------------------------
% ------------------------- USTAWIENIA STRONY--------------------------
% -----------------------------------------------------------------------------------
\usepackage{geometry}
\geometry{
	left=10mm,
	right=10mm,
	top=15mm,
	bottom=15mm,
	marginparwidth=0mm,
	marginparsep=0mm
	}
\newcommand{\wstawplik}[1]{\input{../include/#1}}%skraca kod
\wstawplik{headerofficium}
% elementy zmienne co do headerofficium:
\graphicspath{{../img/}} % katalog obrazów
%\svgpath{{../include/}}
\grechangestaffsize{20}%rozmiar wyniku GregorioTEX
% koniec 
\gresetlastline{ragged}% justowanie do lewej
% debugowanie:
% %\iffalse %*** %\fi %***
% \tracingall przed problematycznym kodem - bardzo dużo informacji
% lualatex+se+debug - dodatkowe opcje: -file-line-error i -halt-on-error 
% \errorcontextlines
%\def\mydebug{1}
\begin{document}
\makeatletter
\grechangedim{interwordspacetext}{\unexpanded{\the\fontdimen2\font plus \the\gluestretch\gre@space@skip@interwordspacenotes minus \the\glueshrink\gre@space@skip@interwordspacenotes}}{fixed}%
\makeatother
\pagestyle{empty}
%\pagestyle{plain}	
\large
\begin{spacing}{0.7}
\begin{center}
	\section*{Cytaty biblijne}
	%\small\textit{Fontevraud Gradual\\}
	\small\textit{Zawsze Wierni nr 4/2001}
	\vspace{4mm}
\end{center}	
\end{spacing}
\setstretch{1}% odstęp między liniami
\small
\selectlanguage{polish}
Fiat lux! – Niech się stanie światłość! (Rodz 1, 3)

Jak na początku Bóg wyprowadził z ciemności światłość, tak też i teraz z niewiernych uczynił nas wiernymi, gdy oświecił nas światłem Wiary. Por. 2 Kor 4, 6: Bóg bowiem, który rozkazał, aby z ciemności zabłysła światłość, sam oświecił serca nasze, abyśmy zajaśnieli blaskiem poznania jasności Bożej w obliczu Chrystusa Jezusa (Quoniam Deus, qui dixit de tenebris lucem splendescere, ipse illuxit in cordibus nostris, ad illuminationem scientiae claritatis Dei, in facie Christi Iesu).
Faciamus hominem ad imaginem et similitudinem nostram – Uczyńmy człowieka na obraz i na podobieństwo nasze (Rodz 1, 26)

Poznajemy w tych słowach tajemnicę Trójcy Przenajświętszej.
Et creavit Deus hominem ad imaginem suam, ad imaginem Dei creavit illum, masculum et feminam creavit eos – I stworzył Bóg człowieka na obraz swój, na obraz Boży stworzył go, mężczyzną i niewiastą stworzył ich (Rodz 1, 27)

Crescite, et multiplicamini, et replete terram, et subicite eam – Rośnijcie i mnóżcie się, i napełniajcie ziemię, a czyńcie ją sobie poddaną (Rodz 1, 28)

Lignum (=Arbor) vitae – Drzewo żywota (Rodz 2, 9)

Było to drzewo ożywiające (arbor vivifica), zachowujące i przedłużające życie, tzn., że spożywający co jakiś czas z niego byliby nieśmiertelni, żyliby w spokoju i radości, wolni od chorób i starości.
Lignumque (Arborque) scientiae boni et mali – I drzewo wiadomości dobrego i złego (Rodz 2, 9)

Przez drzewo to człowiek nauczył się i poznał (na nieszczęście swoje) doświadczalnie (per experientiam) to, co wcześniej znał tylko teoretycznie (per speculationem), mianowicie jaka jest różnica między posłuszeństwem a nieposłuszeństwem, między dobrem a złem, i stąd nazwa: drzewo wiadomości dobrego i złego, tzn. drzewo z którego człowiek poznał przez doświadczenie dobro i zło.
Non est bonum esse hominem solum – Nie jest dobrze być człowiekowi samemu (Rodz 2, 18)

Człowiek jest istotą społeczną (zôn politikon) i stworzony jest do towarzystwa. Vae soli (Biada samotnemu) (Ekle 4, 10).
De fructu ligni, quod est in medio paradisi, praecepit nobis Deus ne comederemus – Z owocu drzewa, które jest w środku raju, rozkazał nam Bóg, abyśmy nie jedli (Rodz 3, 3)

Fructus vetitus (owoc zakazany).
Cumque cognovissent se esse nudos, consuerunt folia ficus – A gdy poznali, że byli nagimi, pozszywali liście figowe (Rodz 3, 7)

Liściem figowym zowiemy niewystarczające okrycie, które zaledwie wstydliwe części ciała zakrywa.
In sudore vultus tui vesceris pane – W pocie oblicza twego będziesz pożywał chleba (Rodz 3, 19)

Jest rozkaz Boży, że chleba z trudem dorabiać się trzeba.
Pulvis es et in pulverem reverteris – Prochem jesteś i w proch się obrócisz (Rodz 3, 19)

Num custos fratris mei sum ego? – Zalim ja jest stróżem brata mego? (Rodz 4, 9)

Vox sanguinis fratris tui clamat ad Me – Głos krwi brata twego woła do Mnie (Rodz 4, 10)

O ciężkim występku mówimy, że woła o pomstę do Boga. Cztery są grzechy, wołające o pomstę: Clamitat ad coelum vox sanguinis et Sodomorum, Vox oppressorum, viduae, pretium famulorum. 1. Rozmyślne zabójstwo. 2. Grzech sodomski. 3. Uciskanie wdów i sierot. 4. Zatrzymywanie zapłaty robotnikom.
Posuitque Dominus Cain signum – I włożył Pan na Kaina znamię (Rodz 4, 15)

O mordercy mówimy, że ma znak Kaina na czole.
Et facti sunt omnes dies Mathusala nongenti sexaginta novem anni – I było wszystkich dni Matusali dziewięćset sześćdziesiąt dziewięć lat (Rodz 5, 27)

Lata Matuzalowe, albo częściej matuzalemowe – przyp. red. Zawsze wierni) znaczy: wiek bardzo długi. Stąd życzenie: obyś dożył lat Matuzalowych!
Quicumque effuderit humanum sanguinem, fundetur sanguis illius – Ktokolwiek by wylał krew człowieczą, będzie wylana krew jego (Rodz 9, 6)

Krew za krew, życie za życie. Prawo odwetu.
Ad imaginem quippe Dei factus est homo – Bo na obraz Boży uczynion jest człowiek (Rodz 9, 6)

Sensus enim et cogitatio humani cordis in malum prona sunt ab adolescentia sua – Zmysł bowiem i myśl serca człowieczego skłonne są do złego od młodości jego (Rodz 8, 21)

Konsekwencje grzechu pierworodnego.
Filii Adam – Synowie Adamowi (Rodz 11, 5)

Filii Adam znaczy: ludzie. To samo znaczy wyrażenie Filius hominis – Syn Człowieczy, którego Chrystus Pan lubił używać z pokory.
Videntes filii Dei filias hominum – Widząc synowie Boży córki ludzkie (Rodz 6,2)

Filii Dei znaczy tu potomków pobożnego Seta. Filius Dei w ścisłym znaczeniu jest Chrystus Pan, w dalszym znaczeniu każdy pobożny człowiek.
Vocatum est nomen eius Babel, quia ibi confusum est labium universae terrae – Nazwano imię jego Babel, gdyż tam pomieszany został język wszystkiej ziemi (Rodz 11, 9)

Babilońskie zamieszanie znaczy tyle, co: wielkie.
Si ad sinistram ieris, ego dexteram tenebo – Jeśli na lewo pójdziesz, ja się udam na prawo (Rodz 13, 9)

Da mihi animas, cetera tolle tibi – Daj mi dusze, inne rzeczy zabierz sobie (Rodz 14, 21)

Dewiza św. Jana Bosco.
Numquid perdes iustum cum impio? – Izali zatracisz sprawiedliwego z niezbożnym? (Rodz 18, 23)

Da mihi de coctione hac rufa – Daj mi z warzywa tego czerwonego (Rodz 25, 30)

Ezaw za misę soczewicy sprzedał bratu prawo pierworodztwa. I my o człowieku, który za lichą cenę zrzeka się ważnego prawa, mówimy, że je sprzedaje za miskę soczewicy.
Det tibi Deus de rore coeli et de pinguedinae terrae – Dajże ci Boże z rosy niebieskiej i z tłustości ziemi (Rodz 27, 28)

Niech ci Bóg da wszystko co najlepsze.
Sint tua tibi – Miej ty swoje (Rodz 33, 9)

Zatrzymaj co masz.
[Castus Ioseph] – por. Rodz 39 (Tak zowiemy młodzieńca czystych obyczajów, który tak nieczystości unika, jak jej unikał Józef Egipski.)

Memento mei, cum bene tibi fuerit – Pamiętaj na mnie, gdy się będziesz miał dobrze (Rodz 40, 14)

Używamy tego cytatu, gdy się polecamy czyjejś pamięci.
Ite ad Ioseph! – Idźcie do Józefa! (Rodz 41, 55)

Hasło czcicieli św. Józefa.
[Filius carissimus Beniamin] – por. Rodz 42, 4 (Najmłodszy i najukochańszy syn Jakuba, stąd najmilsze dziecko, pieszczocha, zowiemy Beniaminkiem.)

Merito haec patimur, quia peccavimus in fratrem nostrum – Słusznie to cierpimy, bośmy zgrzeszyli przeciw bratu naszemu (Rodz 42, 21)

Za grzech zawsze czeka kara.
Vos cogitastis de me malum, sed Deus vertit illud in bonum – Wyście o mnie źle myśleli, ale Bóg obrócił to w dobre (Rodz 50, 20)

Bóg nawet ze zła potrafi wyprowadzić dobro. Przysłowie: nie ma złego, które by na dobre nie wyszło.
[Plagae AEgypti – Plagi egipskie)] (por. Wyjśc 7, 3)

Tyle znaczy, co: wielkie klęski, straszne.
Et factae sunt tenebrae horribiles in universa terra AEgypti – I nastały ciemności straszliwe po wszystkiej ziemi egipskiej (Wyjśc 10, 22)

Egipskie ciemności znaczy tyle, co: wielkie, straszne.
Digitus Dei est hic – Palec to Boży jest (Wyjśc 8, 19)

Gdy stanie się coś nadzwyczajnego lub gdy ciężka kara spadnie, mówimy: to palec Boży.
Ego sum Dominus Deus tuus, qui eduxi te de terra AEgypti, de domo servitutis – Jam jest Pan, Bóg twój, którym cię wywiódł z ziemi egipskiej, z domu niewoli ([1] Non habebis deos alienos coram me – Nie będziesz miał bogów cudzych przede Mną (...) [2] Non assumes nomen Domini Dei tui in vanum – Nie będziesz brał imienia Pana, Boga twego, nadaremno (...) [3] Memento ut diem sabbati sanctifices [- diem dominicum pro die sabbati Apostoli colendum ordinaverunt] – Pamiętaj, abyś dzień sobotni święcił [- dzień niedzielny zamiast sobotniego jest dniem świątecznym z ustanowienia Apostołów (...) [4] Honora patrem tuum et matrem tuam – Czcij ojca twego i matkę twoją (...) [5] Non occides – Nie będziesz zabijał, [6] Non moechaberis – Nie będziesz cudzołożył, [7] Non furtum facies – Nie będziesz kradzieży czynił, [8] Non loqueris contra proximum tuum falsum testimonium – Nie będziesz mówił fałszywego świadectwa przeciw bliźniemu twemu (...) [9] nec desiderabis uxorem eius [sc. proximi tui] – i nie będziesz pragnął żony jego [tj. bliźniego twego]), [10] nec omnia quae illius sunt – ani żadnej rzeczy, która jego jest: Wyjśc 20, 2-17)

Powt Pr 5, 6-20. Dziesięcioro Przykazań Bożych jest treścią wszystkich praw Bożych (Decalogus legum omnium et praeceptorum divinorum est summa et epitome) – św. Augustyn, cyt. przez Katechizm rzymski.
Feceruntque sibi vitulum conflatilem, et adoraverunt – I uczynili sobie cielca odlewanego, i pokłonili się (Wyjśc 32, 8)

Złoty cielec u nas znaczy: złoto, bogactwo. Czcicieli bogactwa zowiemy czcicielami złotego cielca.
(...) oculum pro oculo, dentem pro dente – oko za oko, ząb za ząb (Wyjśc 21, 24)

Prawo odwetu.
Ego enim sum Dominus Deus vester: sancti estote, quia ego sanctus sum – Ja bowiem jestem Pan, Bóg wasz: świętymi bądźcie, bom Ja święty jest (Kapł 11, 44)

[Offeret unum hircum pro peccatis, et alterum dimittet – Będzie ofiarował jednego kozła za grzechy, a drugiego wypuści)] (por. Kapł 16, 18-21)

Na głowę tego drugiego kozła składał arcykapłan wszystkie grzechy ludu, potem go wyganiano na pustynię. My kozłem ofiarnym zowiemy człowieka, na którego inni swe winy zwalają i który za nie pokutować musi.
Testes invoco hodie coelum et terram – Za świadków biorę dziś niebo i ziemię (Powt Pr 4, 26)

Wzywać na świadectwo niebo i ziemię – wszystko.
Audi, Israel! Dominus Deus noster, Dominus unus est – Słuchaj, Izraelu! Pan, Bóg nasz, Pan jeden jest. Hebr.: Szema Israel! Adonai elachenu Adonai echad (Powt Pr 6, 4)

Dominum Deum tuum timebis et illi soli servies – Pana, Boga twego, bać się będziesz i Jemu samemu służyć będziesz (Powt Pr 6, 13, por. Mt 4, 10)

Non tentabis Dominum Deum tuum – Nie będziesz kusił Pana, Boga twego (Powt Pr 6, 16, Mt 4, 7)

Non in solo pane vivit homo, sed in omni verbo quod procedit de ore Dei – Nie samym chlebem żyje człowiek, ale wszelkim słowem, które pochodzi z ust Bożych (Powt Pr 8, 3, Mt 4, 4)

In ore duorum aut trium testium stabit omne verbum – W ustach dwu albo trzech świadków stanie każde słowo (Powt Pr 19, 15)

Pro mensura peccati erit et plagarum modus – Według miary grzechu będzie i razów miara (Powt Pr 25, 2)

Według miary grzechu jest i kary miara. Bo za ciężki grzech ciężka jest i kara.
Tradidit nobis [Deus] terram lacte et melle manantem – Dał nam [Bóg ] ziemię mlekiem i miodem płynącą (Powt Pr 26, 9)

Ziemia, płynąca mlekiem i miodem, znaczy: bardzo urodzajna, mająca dużo bydła i pszczół.
Custodivit [eum Dominus] quasi pupillam oculi sui – Strzegł [go Pan] jako źrenicy oka swego (Powt Pr 32, 10)

Strzec kogo jak źrenicy oka znaczy: strzec bardzo troskliwie.
Homo Dei – Człowiek Boży (Powt Pr 33, 1)

Człowiekiem Bożym zowie Pismo św. Mojżesza, Dawida, proroków – słowem wielkich świętych.
Sicut dies iuventutis tuae, ita et senectus tua – Jako dni młodości twojej, tak i starość twoja (Powt Pr 33, 25)

Jakim kto był za młodu, takim zwykle jest na starość. Przysłowie: czym za młodu skorupka nawre, tym na starość trąci.
Sol, contra Gabaon ne movearis (...) Steteruntque sol – Słońce naprzeciw Gabaonu nie ruszaj się (...) I stanęło słońce (Joz 10, 12-13)

Fakt niezwykły (cud), polegający na przedłużeniu światła słonecznego ponad czas normalny (o 12 lub 24 godziny) – tekst hebr.: jakoby cały dzień. Do tego jednak nie jest rzeczą konieczną, by słońce lub ziemia rzeczywiście zatrzymane zostały w swoim obrocie. Wystarczyło, by nastąpiła refrakcja czyli załamanie promieni słonecznych, które by nadal oświecały widnokrąg ziemski w tym samym miejscu, chociaż słońce posuwało się naprzód ku zachodowi lub też już zupełnie zaszło, na podobieństwo zjawiska naturalnego, stwierdzonego przez naukę przy wschodzie i zachodzie słońca. Dzięki bowiem temu załamaniu się promieni słonecznych widzimy rano słońce parę minut wcześniej na horyzoncie niż ono zjawi się rzeczywiście czyli według obliczeń astronomicznych, a przy zachodzie słońca widzimy je parę minut dłużej już po rzeczywistym zniknięciu z horyzontu. Otóż Pan Bóg to naturalne zjawisko w sposób cudowny mógł przedłużyć na przeciąg całego dnia, nie wstrzymując wcale naturalnego obrotu ciał niebieskich.
Ego autem et domus mea serviemus Domino – A ja i dom mój Panu służyć będziemy (Joz 24, 15)

Sicut feci [aliis], ita reddidit mihi Deus – Jakom czynił [innym], tak mi Bóg oddał (Sędz 1, 7)

Jak ty bliźniemu, tak Bóg tobie.
Si non arassetis in vitula mea, non invenissetis propositionem meam – Byście byli nie orali jałowicą moją, nie rozwiązalibyście zagadki mojej (Sędz 14, 18)

Tak rzekł Samson do Filistynów, gdy z pomocą jego żony odgadli zadaną im zagadkę.
Philisthiim super te, Samson! – Filistyni nad tobą, Samsonie! (Sędz 16, 9)

Pax tecum sit! – Pokój z tobą niech będzie! (Sędz 19, 20)

Haec mihi faciat Dominus et haec addat – To mi niech uczyni Pan i to niech przyczyni (Rut 1, 17)

Takiej formuły używali żydzi, zaklinając się lub przysięgając.
Deus scientiarum Dominus est – Bóg umiejętności Pan jest (1 Król 2, 3)

Quicumque glorificaverit me, glorificabo eum – Ktokolwiek Mię wielbić będzie, uwielbię go (1 Król 2, 30)

Ecce ego, quia vocasti me – Otom ja, boś mię wołał! (1 Król 3, 6)

Samuel do Helego. Gotowość do pełnienia rozkazów Bożych.
Loquere, Domine, quia audit servus tuus – Mów, Panie, bo słucha sługa Twój (1 Król 3, 9)

Num et Saul inter prophetas? – Czyż i Saul między prorokami? (1 Król 10, 11)

Coś niesłychanego.
Quaesivit Dominus sibi virum iuxta cor suum – Poszukał sobie Pan męża według serca Swego (1 Król 13, 14)

Dawida nazywa Pismo św. mężem według serca Bożego.
Pro eo ergo quod abiecisti sermonem Domini, abiecit te Deus – Za to tedy, żeś odrzucił mowę Pańską, odrzucił cię Pan (1 Król 15, 23)

Siccine separat amara mors? – Tak to rozłącza gorzka śmierć? (1 Król 15, 32)

Vir (...) nomine Goliath (...), altitudinis sex cubitorum – Mąż (...) imieniem Goliat (...), wysoki na sześć łokci (1 Król 17, 4)

Męża rosłego i silnego zowiemy Goliatem.
Vivit Dominus et vivit anima tua – Żyje Pan i żyje dusza twoja (1 Król 20, 3)

Formuła przysięgi żydowskiej.
Uno tantum (ut ita dicam) gradu ego morsque dividimur – Jeden tylko, że tak rzekę, krok od śmierci mnie dzieli (1 Król 20, 3)

Byłem krok od śmierci mówimy, gdyśmy uniknęli groźnego niebezpieczeństwa.
Si inveni gratiam in oculis tuis – Jeślim znalazł łaskę w oczach twoich (1 Król 20, 29)

Te słowa znaczą: jeżeli łaskę raczysz mi okazać.
Iudicet Dominus inter me et te – Niech rozsądzi Pan między mną a tobą (1 Król 24, 13)

Niech Bóg rozsądzi między nami, który z nas ma słuszność.
Ab impiis egredietur impietas – Od niezbożnych wynijdzie niezbożność (1 Król 24, 14)

Jakie drzewo, taki owoc, z ciernia nie będzie winogron.
Tu enim tribuisti mihi bona, ego autem reddidi tibi mala – Boś ty mnie wyświadczył dobro, a jam tobie złem się odpłacił (1 Król 24, 18)

Rectus es tu et bonus in conspectu meo (...) sed satrapis non places – Tyś prawy i dobry w oczach moich (...) ale książętom się nie podobasz (1 Król 29, 6)

Aequa enim pars erit descendentibus ad praelium et remanentis ad sarcinas – Równy bowiem dział będzie idącego do bitwy i zostającego przy tłumokach (1 Król 30, 24)

An ignoras quod periculosa sit desperatio? – Czyż nie wiesz, że bardzo niebezpieczna jest rozpacz? (2 Król 2, 26)

Do rozpaczy kogo doprowadzić – rzecz niebezpieczna.
Manete in Iericho donec crescat barba vestra – Mieszkajcie w Jerychu, aż urośnie broda wasza (2 Król 10, 5)

Filius mortis est vir qui fecit hoc! – Synem śmierci jest mąż, który to uczynił! (2 Król 12, 5)

Dawid do Natana.
Tu es ille vir! – Ty jesteś tym człowiekiem! (2 Król 12, 7)

Natan do Dawida.
Peccavi Domino! – Zgrzeszyłem Panu! (2 Król 12, 13)

Dawid do Natana.
Cum sancto sanctus eris, (...) et cum perverso perverteris – Ze świętym święty będziesz, (...) a z przewrotnym przewrotny się staniesz (2 Król 22, 26-27)

Z jakim przestajesz, takim się stajesz.
Ego ingredior viam universae terrae – Ja idę w drogę wszystkiej ziemi (3 Król 2, 2)

To jest: ja zdążam do śmierci, jak pozostali ludzie zamieszkujący ziemię, tak, iż z prochu (ziemi) w proch (ziemię) się obrócę. Por. Rodz 3, 19.
[Iudicium Salomonis – Sąd Salomona lub wyrok salomonowy)] (por. 3 Król 3, 16-28)

Gdy ktoś nadzwyczaj trudny wypadek dobrze i sprawiedliwie osądzi, to mówimy, że sądzi jak Salomon. Znany jest sąd Salomona w sprawie dwóch matek i ich dwojga dzieci.
Dedit quoque Deus sapientiam Salomoni et prudentiam multam nimis (...) Et praecedebat sapientia Salomonis sapientiam omnium Orientalium et AEgyptiorum, et erat sapientior cunctis hominibus – Dał też Bóg Salomonowi mądrość i roztropność bardzo wielką (...) I przechodziła mądrość Salomona mądrość wszystkich ludzi wschodnich i Egipcjan. I był mędrszy od wszystkich ludzi (3 Król 4, 29-31)

O człowieku bardzo mądrym mówimy: mądry jak Salomon.
Dormivitque Salomon cum patribus suis – I zasnął Salomon z ojcami swymi (3 Król 11, 43)

Znaczy tyle, co: umarł.
Usquequo claudicatis in duas partes? Si Dominus est Deus, sequimini eum, si autem Baal, sequimini illum – Pókiż będziecie chromać na dwie strony? Jeśli Pan jest Bogiem, idźcie za Nim, a jeśli Baal, idźcie za nim! (3 Król 18, 21)

Nie chwiejcie się na dwie strony, lecz trzymajcie się prawdy.
Et derelinquam mihi in Israel septem millia virorum, quorum genua non sunt incurvata ante Baal – I zostawię sobie w Izraelu siedem tysięcy mężów, których kolana nie zginały się przed Baalem (3 Król 19, 18)

Pan Bóg ma zawsze dusze wybrane.
Ego odi eum, quia non prophetat mihi bonum, sed malum – Ja go nienawidzę, ponieważ mi nie prorokuje dobrego, ale zło (3 Król 22, 8)

Przepowiadaczy nieszczęść zwykle niecierpimy i prześladujemy – prawda w oczy kole.
Haeccine est illa Iezabel? – Taż to jest owa Jezabel? (4 Król 9, 37)

Zaiste można tak było zawołać, gdy potężną niegdyś królową psy szarpały.
Servi tui sumus – quaecumque iusseris faciemus (Jesteśmy służebnicy twoi, cokolwiek rozkażesz, uczynimy)

4 Król 10, 5
Nolite tangere christos meos – Nie tykajcie pomazańców Moich (1 Kron 16, 22)

Nie porywajcie się na wybrańców Moich.
Melius mihi est, ut incidam in manus Domini, quia multae sunt miserationes eius, quam in manus hominum – Lepiej mi wpaść w ręce Pańskie, bo mnogie są miłosierdzia Jego, niźli w ręce człowiecze (1 Kron 21, 13)

Si quaesieris eum (sc. Deum), invenies, si autem dereliqueris eum, proiciet te in aeternum – Jeśli Go [tj. Boga] szukać będziesz, znajdziesz, ale jeśli Go opuścisz, odrzuci cię na wieki (1 Kron 28, 9)

Peregrini enim sumus coram te, [Deus], et advenae, sicut omnes patres nostri – Bośmy goście przed Tobą, [Boże], i przechodnie, jako i wszyscy ojcowie nasi (1 Kron 29, 15)

Tu enim solus [Domine] nosti corda filiorum hominum – Ty bowiem sam [jeden, Panie,] znasz serca synów ludzkich (2 Kron 6, 30)

Oculi enim Domini contemplantur universam terram, et praebent fortitudinem his qui corde perfecto credunt in eum – Albowiem oczy Pańskie patrzą na wszystką ziemię i dodają mocy tym, którzy sercem doskonałym weń wierzą (2 Kron 16, 9)

Stulte igitur egisti, et propter hoc ex praesenti tempore adversum te bella consurgent – A tak głupioś uczynił i dlatego od tego czasu powstaną przeciwko tobie walki (2 Kron 16, 9)

Non est enim apud Dominum Deum nostrum iniquitas, nec personarum acceptio – Bo nie masz u Pana, Boga naszego, nieprawości ani względu na osoby (2 Kron 19, 7)

Et convertit Deus noster maledictionem in benedictionem – I obrócił Bóg nasz przekleństwo w błogosławieństwo (Neh 13, 2)

Et nunc, Domine, memor esto mei, et ne vindictam sumas de peccatis meis, neque reminiscaris delicta mea, vel parentum meorum – A teraz, Panie, wspomnij na mnie, nie mścij się za grzechy moje ani pamiętaj na występki moje albo rodziców moich (Tob 3, 3)

Quod ab alio oderis fieri tibi, vide ne tu aliquando alteri facias – Czego byś nie chciał, aby ci kto inny czynił, patrz, abyś ty kiedy drugiemu nie czynił (Tob 4, 16)

Consilium semper a sapiente perquire – Rady zawsze u mądrego szukaj (Tob 4, 19)

Qui autem faciunt peccatum et iniquitatem, hostes sunt animae suae – A ci, którzy czynią grzech i nieprawość, są nieprzyjaciółmi duszy swojej (Tob 12, 10)

Nuntius venit ad Iob (...) Cumque adhuc ille loqueretur, venit alter... – Przybiegł do Joba poseł (...) A gdy ten jeszcze mówił, przyszedł drugi... (Job 1, 14-19)

Wieść Hiobowa znaczy tyle, co: niespodziewana a straszna. Tak samo poselstwo Jobowe znaczy tyle, co: zwiastujące nieszczęście.
Dominus dedit, Dominus abstulit, sicut Domino placuit, ita factum est. Sit nomen Domini benedictum – Pan dał, Pan wziął: jako się Panu podobało, tak się stało. Niech będzie imię Pańskie błogosławione! (Job 1, 21)

Cuncta quae habet homo dabit pro anima sua – Wszystko, co ma człowiek, da za duszę swoją (Job 2, 4)

Quasi una de stultis mulieribus locuta es – Jako jedna z niewiast głupich rzekłaś (Job 2, 10)

Żona Hioba uważała, że skoro Bóg zsyła cierpienia na niewinnego, to w takim razie nie warto Mu służyć, nie warto być dobrym. Hiob słusznia oburza się na takie poglądy.
Si bona suscepimus de manu Dei, mala quare non suscipiamus? – Jeśliśmy przyjęli dobra z ręki Bożej, złego czemu byśmy przyjmować nie mieli? (Job 2, 10)

Vere stultum interficit iracundia – Zaiste gniew zabija głupiego (Job 5, 2)

Nihil in terra sine causa fit – Nic się na ziemi bez przyczyny nie dzieje (Job 5, 6)

Militia est vita hominis super terram – Bojowaniem jest żywot człowieczy na ziemi (Job 7, 1)

Hesterni quippe sumus, et ignoramus, quoniam sicut umbra dies nostri sunt super terram – Wczorajsi bowiem jesteśmy i nic nie wiemy, gdyż dni nasze na ziemi są jako cień (Job 8, 9)

Deridetur enim iusti simplicitas – Naśmiewają się bowiem z prostoty sprawiedliwego (Job 12, 4)

Etiam si occiderit me, in ipso sperabo – Choćby mię zabił, w Nim ufać będę (Job 13, 15)

Słowa te wyrażają bezgraniczne zaufanie do Boga.
Homo, natus de muliere, brevi vivens tempore, repletur multis miseriis – Człowiek, urodzony z niewiasty, żyjąc przez czas krótki, napełnion bywa wielu nędzami (Job 14, 1)

Miseremini mei, miseremini mei, saltem vos, amici mei – Zmiłujcie się nade mną, zmiłujcie się nade mną, chociaż wy, przyjaciele moi (Job 19, 21)

Hoc scio a principio, ex quo positus est homo super terram, quod laus impiorum brevis sit, et gaudium hypocritae ad instar puncti – To wiem od początku, odtąd, jako osadzon jest człowiek na ziemi, że chwała niezbożnych krótka jest, a radość obłudnika jako mgnienie oka (Job 20, 4-5)

Quid prodest Deo, si iustus fueris? – Co za pożytek Bogu jest, jeśli sprawiedliwym będziesz? (Job 22, 3)

Zob. komentarz św. Grzegorza Wielkiego: We wszystkim, co dobrze czynimy, nie Boga, ale nas samych wspieramy. In omni quod bene agimus nosmetipsos, non Deum iuvamus – Moral. XVI, I
Si peccaveris, quid ei (sc. Deo) nocebis? – Jeśli zgrzeszysz, cóż Mu (tj. Bogu) zaszkodzisz? (Job 35, 6)

Kto grzeszy, ten przede wszystkim szkodzi sobie samemu.
Pepigi foedus cum oculis meis, ut ne cogitarem quidem de virgine – Uczyniłem przymierze z oczyma swymi, abym ani pomyślał o pannie (Job 31, 1)

Postanowienie unikania okazji do grzechu. Najwyższy stopień czystości przejawia się w obrzydzeniu do wszelkich myśli nieskromnych (św. Bonawentura).
Beatus vir qui non abiit in consilio impiorum, et in via peccatorum non stetit, et in cathedra pestilentiae non sedit – Błogosławiony mąż, który nie chodził za radą niezbożnych i na drodze grzesznych nie stał, i na stolicy zaraźliwości nie siedział (Ps 1, 1)

Szczęśliwy ten człowiek, który nie idzie za radami, prawami i opiniami ludzi bezbożnych.
Quoniam novit Dominus viam iustorum, et iter impiorum peribit – Albowiem zna Pan drogę sprawiedliwych, a droga niezbożnych zginie (Ps 1, 6)

Bóg zna, tj. pochwala postępowanie sprawiedliwych, lecz postępowanie grzeszników prowadzi na zatracenie.
Astiterunt reges terrae, et principes convenerunt in unum adversus Dominum, et adversus Christum eius – Stanęli wespół królowie ziemscy, a książęta zeszli się w gromadę przeciw Panu i przeciw Chrystusowi Jego (Ps 2, 2)

Psalmista w duchu proroczym zapowiada bunt i walkę, jaką stale Chrystusowi Panu i Jego Kościołowi będą wypowiadać królowie i rządy owiane duchem doczesności.
Servite Domino in timore, et exultate ei cum tremore – Służcie Panu w bojaźni, a radujcie się Mu ze drżeniem (Ps 2, 11)

Bogu należy oddawać cześć z wielkim uszanowaniem.
Irascimini, et nolite peccare – Gniewajcie się, ale nie grzeszcie (Ps 4, 5. Por. Ef 4, 26)

Jeżeli chcecie krytykować, to i owszem, ale zanim coś powiecie wpierw nad tym dobrze się zastanówcie.
In pace (...) dormiam, et requiescam – W pokoju będę (...) spał i odpoczywał (Ps 4, 9)

My mówimy w modlitwach Kościoła o zmarłym: Requiescat in pace!
Dixit insipiens in corde suo: Non est Deus – Rzekł głupi w sercu swoim: Nie ma Boga (Ps 13, 1)

Nolite fieri sicut equus et mulus, quibus non est intellectus – Nie bądźcież jako koń i muł, które nie mają rozumu (Ps 31, 9)

Abyssus abyssum invocat – Przepaść przepaści przyzywa (Ps 41, 8)

Nieskończenie doskonały Bóg i otchłań naszej nędzy.
Miserere mei, Deus, secundum magnam misericordiam tuam, et secundum multitudinem miserationum tuarum, dele iniquitatem meam – Zmiłuj się nade mną, Boże, według wielkiego miłosierdzia Twego, i według mnóstwa litości Twoich zgładź nieprawość moją! (Ps 50, 1)

Cor contritum et humiliatum, Deus, non despicies – Sercem skruszonym i uniżonym, Boże, nie wzgardzisz! (Ps 50, 19)

Tu [Domine] reddes unicuique iuxta opera sua – Ty [Panie] oddasz każdemu według uczynków jego (Ps 61, 13)

In te, Domine, speravi, non confundar in aeternum – W Tobiem, Panie, nadzieję miał, niech nie będę pohańbion na wieki! (Ps 70, 1)

Misericordias Domini in aeternum cantabo, in generationem et generationem annuntiabo veritatem tuam in ore meo – Zmiłowania Pańskie na wieki wyśpiewywać będę, od pokolenia do pokolenia prawdę Twoją usty mymi opowiadać będę (Ps 88, 2)

Quoniam angelis suis mandavit de te, ut custodiant te in omnibus viis tuis. In manibus portabunt te, ne forte offendas ad lapidem pedem tuum – Albowiem aniołom twoim rozkazał o tobie, aby cię strzegli na wszystkich drogach twoich. Na ręku będą cię nosić, byś snadź nie uraził o kamień nogi twojej (Ps 90, 11-12)

Na rękach kogo nosić znaczy: bardzo troskliwie się nim opiekować.
Non nobis, Domine, non nobis, sed nomini tuo da gloriam – Nie nam, Panie, nie nam, ale imieniu Twemu daj chwałę (Ps 113B, 1)

Omnis homo mendax – Każdy człowiek kłamca! (Ps 115, 11)

Kłamca: zawodny, jedynie Bogu można bezpiecznie zaufać.
Lapidem quem reprobaverunt aedificantes, hic factus est in caput anguli – Kamień, który odrzucili budujący, ten się stał kamieniem węgielnym (Ps 117, 22)

Wyrażenie przysłowiowe: to, co zostało odrzucone jako nic nie warte, okazało się najlepszym i najważniejszym. Takim jest Chrystus Pan, odrzucony przez przywódców narodu żydowskiego.
Haec est dies quam fecit Dominus – Ten jest dzień, który uczynił Pan (Ps 117, 24)

Benedictus qui venit in nomine Domini – Błogosławiony, który przychodzi w imię Pańskie! (Ps 117, 26)

Qui seminant in lacrymis, in exultatione metent – Którzy sieją ze łzami, będą żąć z radością (Ps 125, 5)

De profundis clamavi ad te, Domine, Domine, exaudi vocem meam – Z głębokości wołałem ku Tobie, Panie, Panie, wysłuchaj głosu mego! (Ps 129, 1-2)

Timor Domini principium sapientiae. Sapientiam atque doctrinam stulti despiciunt – Bojaźń Pańska początkiem mądrości, mądrością i nauką głupi gardzą (Przyp 1, 7)

Disciplinam Domini, fili mi, ne abicias, nec deficias cum ab eo corriperis, quem enim diligit Deus corripit, et quasi pater in filio complacet sibi – Karania Pańskiego, synu mój, nie odrzucaj ani ustawaj, gdy od Niego karan bywasz, bo kogo Pan miłuje, karze i jako ojciec ma upodobanie w synu (Przyp 3, 11-12)

Beatus homo qui invenit sapientiam, et qui affluit prudentia – Błogosławiony człowiek, który znalazł mądrość i który obfituje w roztropność (Przyp 3, 13)

Noli prohibere benefacere eum qui potest, si vales, et ipse benefac – Nie zabraniaj dobrze czynić temu, który może, jeśli możesz, i sam dobrze czyń (Przyp 3, 27)

Ne moliaris amico tuo malum, cum ille in te habeat fiduciam – Nie obmyślaj przeciw przyjacielowi twemu zła, gdy on w tobie ufność pokłada (Przyp 3, 29)

Gloriam sapientes possidebunt, stultorum exaltatio ignominia – Chwałę mądrzy osiągną, głupich wywyższeniem zelżywość (Przyp 3, 35)

Tene disciplinam, ne dimittas eam, custodi illam, quia ipsa est vita tua – Trzymaj się nauki, nie opuszczaj jej, strzeż jej, bo ona jest życiem twoim (Przyp 4, 13)

Ne delecteris in semitis impiorum, nec tibi placeat malorum via – Nie kochaj się w ścieżkach bezbożników ani sobie podobaj w drodze złośników (Przyp 4, 14)

Via impiorum tenebrosa, nesciunt ubi corruant – Droga bezbożnych ciemna, nie wiedzą, gdzie upadają (Przyp 4, 19)

Respicit Dominus vias hominis, et omnes gressus eius considerat – Pan patrzy na drogi człowieka, i uważa na wszystkie kroki jego (Przyp 5, 21)

Homo apostata vir inutilis, graditur ore perverso (...) Huic extemplo veniet perditio sua, et subito conteretur, nec habebit ultra medicinam – Człowiek odstępca, mąż nieużyteczny, chodzi z przewrotnymi ustami (...) Temu nagle przyjdzie zatracenie jego i niespodzianie skruszony [starty] będzie, a nie będzie miał więcej lekarstwa (Przyp 6, 12-15)

Conserva, fili mi, praecepta patris tui, et ne dimittas legem matris tuae – Zachowaj, synu mój, przykazania Ojca twego, a nie opuszczaj wskazań Matki twojej (Przyp 6, 20, por. 1, 8)

Obowiązek przestrzegania przykazań Bożych i kościelnych.
Veritatem meditabitur guttur meum, et labia mea detestabuntur impium – Prawdę rozmyślać będzie gardło moje, a wargi moje będą się brzydzić bezbożnością (Przyp 8, 7)

Melior est enim sapientia cunctis pretiosissimis, et omne desiderabile ei non potest comparari – Lepsza jest bowiem mądrość nad wszelakie najdroższe rzeczy, i wszystko, co może być pożądane, nie może być z nią porównane (Przyp 8, 11)

Ubi fuerit superbia, ibi erit et contumelia. Ubi autem est humilitas, ibi est sapientia – Gdzie będzie pycha, tam będzie i hańba, a gdzie jest pokora, tam i mądrość (Przyp 11, 2)

Doctrina bona dabit gratiam, in itinere contemptorum vorago – Nauka dobra da łaskę, na drodze gardzących przepaść (Przyp 13, 15)

Cor sapientis quaerit doctrinam, et os stultorum pascitur imperitia – Serce mądrego szuka nauki, a usta głupich karmią się nieumiejętnością (Przyp 15, 14)

Melius est nomen bonum quam divitiae multae – Lepsze jest imię dobre niż wielkie bogactwa (Przyp 22, 1)

Ne transgrediaris terminos antiquos quos posuerunt patres tui – Nie przestępuj granic starych, które założyli ojcowie twoi (Przyp 22, 28)

Vanitas vanitatum, et omnia vanitas – Marność nad marnościami i wszystko marność (Ekle 1, 2)

Autor De imitatione Christi dodaje: praeter amare Deum et illi soli servire – oprócz [jednego:] miłować Boga i służyć Jemu samemu, lib. I, cap. 1, 3
Melior est canis vivus leone mortuo – Lepszy jest pies żywy aniżeli lew zdechły (Ekle 9, 4)

Fortis est ut mors dilectio – Mocna jest jak śmierć miłość (PnP 8, 6)

Initium sapientiae timor Domini – Początkiem mądrości bojaźń Pańska (Ekli 1, 16)

Initium omnis peccati est superbia – Pycha jest początkiem grzechu każdego (Ekli 10, 15)

Qui tetigerit picem inquinabitur ab ea, et qui communicaverit superbo induet superbiam – Kto się smoły dotknie, pomaże się nią, a kto ma towarzystwo z pysznym, oblecze się w pychę (Ekli 13, 1)

Mitte illum in operationem, ne vacet, multam enim malitiam docuit otiositas – Poślij go na robotę, aby nie próżnował, bo wielkiej złości nauczyło próżnowanie (Ekli 33, 28-29)

Omnia opera Domini bona – Wszystkie dzieła Pańskie są dobre (Ekli 39, 39)

Quiescite agere perverse, discite benefacere – Przestańcie źle czynić, uczcie się czynić dobrze (Iz 1, 16-17)

Vae qui dicitis malum bonum, et bonum malum, ponentes tenebras lucem, et lucem tenebras, ponentes amarum in dulce, et dulce in amarum! – Biada wam, którzy nazywacie złe dobrym, a dobre złym, zamieniacie ciemność na światłość, a światłość na ciemność, zamieniacie gorzkie na słodkie, a słodkie na gorzkie! (Iz 5, 20)

Vae qui sapientes estis in oculis vestris, et coram vobismetipsis prudentes! – Biada wam, którzy jesteście mądrymi w oczach waszych i sami przed sobą roztropnymi! (Iz 5, 21)

Vae qui potentes estis ad bibendum vinum, et viri fortes ad miscendam ebrietatem! – Biada wam, którzy jesteście mocnymi w piciu wina, i mężami dzielnymi w mieszaniu opilstwa! (Iz 5, 22)

Vae mihi, quia tacui! – Biada mnie, że milczałem! (Iz 6, 5)

Comedamus et bibamus, cras enim moriemur – Jedzmy i pijmy, bo jutro pomrzemy! (Iz 22, 13, 1 Kor 15, 32)

Hasło ludzi cielesnych.
Populus hic labiis me honorat, cor autem eorum longe est a me. Sine causa autem colunt me, docentes doctrinas et mandata hominum – Lud ten czci Mię wargami, ale serce ich daleko jest ode Mnie, lecz na próżno cześć Mi oddają, gdy głoszą nauki i ustawy ludzkie (Iz 29, 13, Mt 15, 8-9)

Stultus enim fatua (=stulta) loquetur – Bo głupi głupie rzeczy mówić będzie (Iz 29, 6)

Omnis caro foenum, et omnis gloria eius quasi flos agri – Wszelkie ciało trawa, a wszelka chwała jego jako kwiat polny (Iz 40, 6)

Non est pax impiis, dicit Dominus – Nie masz pokoju niezbożnym, mówi Pan (Iz 48, 22)

Et curabant (...) dicentes: Pax, pax! et non erat pax – I leczyli (...) mówiąc: Pokój, pokój! – a nie było pokoju (Jer 6, 14)

Patres comederunt uvam acerbam, et dentes filiorum obstupuerunt – Ojcowie jedli jagodę kwaśną, a zęby synów ścierpły (Jer 31, 29, Ez 18, 2)

Za grzechy ojców cierpią nieraz dzieci.
Grex perditus factus est populus meus, pastores eorum seduxerunt eos – Trzodą straconą stał się lud mój, pasterze ich zwiedli ich (Jer 50, 6)

Por. przysłowie: Biada wtedy trzodzie, gdy wilcy z pasterzami w zgodzie.
Curavimus Babylonem, et non est sanata, derelinquamus eam – Leczyliśmy Babilon, a nie jest uzdrowiony. Opuśćmy go! (Jer 51, 9)

Quomodo sedet sola civitas plena populo! Facta est quasi vidua domina gentium, princeps provinciarum facta est sub tributo – Jakże samotnie siedzi miasto, pełne ludu! Stała się jakby wdową pani narodów, księżna powiatów stała się hołdowną (Lam 1, 1)

Anima quae peccaverit, ipsa morietur – Dusza, która zgrzeszy, ta umrze (Ez 18, 4)

Vae pastoribus (...), qui pascebant semetipsos! Nonne greges a pastoribus pascuntur? – Biada pasterzom (...), którzy paśli samych siebie! Izali nie trzody pasą pasterze? (Ez 34, 2)

Mane, Thecel, Phares – Mane, Tekel, Fares (Dan 5, 25)

Policzono, zważono, rozdzielono.
Appensus es in statera, et inventus es minus habens – Zważony jesteś na wadze i znaleziono cię za lekkim (Dan 5, 27)

Et polluent sanctuarium fortitudinis, et auferent iuge sacrificium, et dabunt abominationem in desolationem. Et impii in testamentum simulabunt fraudulenter – I zgwałcą świątynię mocy, i odejmą ustawiczną ofiarę, a dadzą obrzydłość na spustoszenie. A bezbożnie działający przeciw przymierzu będą postępować obłudnie i zdradliwie (Dan 11, 31-32)

Propter hoc ecce ego lactabo eam, et ducam eam in solitudinem, et loquar ad cor eius – Przeto oto Ja przynęcę ją i zawiodę ją na puszczę, i będę mówił do serca jej (Oz 2, 14)

Bóg powołuje niektóre dusze do życia kontemplacyjnego w samotności.
Quia tu scientiam repulisti, repellam te – Ponieważ ty odrzuciłeś umiejętność, odrzucę cię (Oz 4, 6)

Przestroga szczególnie dla kapłanów. Por. Mal 2, 7.
Qui nesciunt quid sit inter dexteram et sinistram suam – Którzy nie rozróżniają między prawicą a lewicą swoją (Jon 4, 11)

Znaczy tyle, co: małe niemowlęta.
Domine (...) cum iratus fueris, misericordiae recordaberis – Panie (...) gdy się rozgniewasz, wspomnisz na miłosierdzie (Hab 3, 2, por. Tob 3, 13)

Non sum propheta, homo agricola ego sum, quoniam Adam exemplum meum ab adolescentia mea – Nie jestem prorokiem, człowiek oracz ja jestem, bo Adam przykładem moim od młodości mojej (Zach 13, 5)

Percute pastorem, et dispergentur oves – Uderz pasterza, a rozproszą się owce (Zach 13, 7, Mt 26, 31)

Sicut autem fuerit voluntas in coelo, sic fiat! – A jaka będzie wola na niebie, tak niechaj się stanie! (1 Mach 3, 60)

Ipsi autem non erant de semine virorum illorum per quos salus facta est in Israel – A oni nie byli z nasienia onych mężów, przez których się stało zbawienie w Izraelu (1 Mach 5, 62)

Verum non propter locum gentem, sed propter gentem locum Deus elegit – Lecz nie dla miejsca naród, ale dla narodu miejsce Bóg obrał (2 Mach 5, 19)
\end{document}